% GENERATED FILE. Edit canonical lessons, then run npm run generate:books.
% canonical-source-hash: fnv1a64:bcb00dc801aa9fb4
% canonical-lessons: ZH-C302-huaiyun, ZH-C302-huo, ZH-W302-huo-alive, ZH-C302-yinger, ZH-C302-jiuhuo-revive, ZH-C302-shuangbaotai

\chapter{A Baby on the Way}
\label{ch:zh-a-baby-on-the-way}
\hlchaptermodality{302}

\begin{chapteropening}
\textbf{By the end of this chapter:} \emph{I can talk about a pregnancy, a newborn, twins, and a life that was saved.}

Complete the last lesson of chapter 302: I can talk about a pregnancy, a newborn, twins, and a life that was saved.
\end{chapteropening}

\section[huáiyùn]{huáiyùn — to be pregnant; to be expecting a baby}

\label{lesson:ZH-C302-huaiyun}



\begin{quote}
Before the new one: say the Mandarin for to break a bone, then the Mandarin for sure enough.
\end{quote}

\subsection*{You'll want to know: huáiyùn}
\emph{huáiyùn} — \textquotedblleft{}to be pregnant; to be expecting a baby\textquotedblright{}.

\begin{quote}
\emph{tā huáiyùn le} — She's pregnant.
\end{quote}

Its characters come later; for now it is a word for the ear.

\begin{cousinweb}
\emph{huáiyùn} is built from \emph{huái} 'to hold close' and \emph{yùn} 'to carry a child', both inherited from Middle Chinese, \emph{hweaj} and \emph{yingH} (Guangyun, in Baxter's transcription), which glosses the first 'to embrace' and the second with this very word. \emph{yingH} would regularly give yìng; the final of yùn is irregular.
\end{cousinweb}

\begin{grammarlens}[title={What you've built}]
One more word for telling what happened.
\end{grammarlens}

\subsection*{Your turn}
\begin{itemize}
\raggedright
  \item \emph{Say it:} \emph{huáiyùn}
  \item \emph{Say it:} \emph{huáiyùn}, once more
  \item \emph{Say it:} \emph{tā huáiyùn le} — She's pregnant.
  \item \emph{Recall:} say the Mandarin for to break a bone, then the Mandarin for sure enough, then say \emph{huáiyùn} again
\end{itemize}

\begin{tcolorbox}[breakable,colback=teal!4,colframe=teal!35!black,title={Before you move on}]
What does huáiyùn mean? (\textquotedblleft{}To be pregnant; to be expecting a baby\textquotedblright{}.) Say it once more.
\end{tcolorbox}
\section[\texorpdfstring{\zh{活} (huó)}{活 (huó)}]{\zh{活} (huó) — to live; alive}

\label{lesson:ZH-C302-huo}



\begin{quote}
Before the new one: say the Mandarin for sure enough, then the Mandarin for to be expecting a baby.
\end{quote}

\subsection*{You'll want to know: \zh{活}}
\textbf{\zh{活}} — \emph{huó} — \textquotedblleft{}to live; alive\textquotedblright{}.

\begin{quote}
\textbf{\zh{他活了一百岁}} \emph{tā huó le yìbǎi suì} — He lived to a hundred.
\end{quote}

The new character, in pieces: \zh{活} — nine strokes in three pieces: the water piece of \zh{洗} on the left, three strokes; then \zh{千}, three strokes; then \zh{口} under it. The next lesson puts it in your hand.

\begin{cousinweb}
\textbf{\zh{活}} (\emph{huó}) is inherited from Middle Chinese \emph{hwat} (Guangyun, in Baxter's transcription), which glosses it 'not dead' and also 'the sound of flowing water'.
\end{cousinweb}

\begin{grammarlens}[title={What you've built}]
One more word for telling what happened.
\end{grammarlens}

\subsection*{Your turn}
\begin{itemize}
\raggedright
  \item \emph{Say it:} \emph{huó}
  \item \emph{Say it:} \emph{huó}, once more
  \item \emph{Say it:} \emph{tā huó le yìbǎi suì} — He lived to a hundred.
  \item \emph{Recall:} say the Mandarin for sure enough, then the Mandarin for to be expecting a baby, then say \emph{huó} again
\end{itemize}

\begin{tcolorbox}[breakable,colback=teal!4,colframe=teal!35!black,title={Before you move on}]
What does \zh{活} mean? (\textquotedblleft{}To live; alive\textquotedblright{}.) Say it once more.
\end{tcolorbox}
\section[\texorpdfstring{\zh{活} (huó)}{活 (huó)}]{\zh{活} — the water piece, \zh{千}, then \zh{口}}

\label{lesson:ZH-W302-huo-alive}



\begin{quote}
The word first, then the character inside it.

\begin{itemize}
\raggedright
  \item \emph{Recall:} say \emph{huó} — \textbf{R1}, one lesson back, before it decays
\end{itemize}

\begin{quote}
\zh{洗} \zh{千}
\end{quote}

\emph{Read it:} each one aloud, with nothing said first — \emph{xǐ}, \emph{qiān}
\end{quote}

\subsection*{Script — the water piece, \zh{千}, then \zh{口}}
\hlblockfigure{\detokenize{figures/ZH-W302-huo-alive-filmstrip.pdf}}{How \zh{活} is written, stroke by stroke}

\begin{quote}
\zh{活}
\end{quote}

\textbf{Three pieces: the water piece on the left, then \zh{千} over \zh{口} on the right.}

\textbf{The water piece, three strokes.} As on the left of \zh{洗}: two dots, then a stroke rising to the right.

\textbf{\zh{千}, three strokes.} A short sweep down to the left, a bar across, and a vertical down through the bar.

\textbf{\zh{口}, three strokes}, under \zh{千}: a short vertical, a stroke that runs right and turns down, and the bar that closes the foot.

3 + 3 + 3 = nine.

\subsection*{Your turn}
\begin{itemize}
\raggedright
  \item \emph{Write it:} \zh{活}, all nine, with the model in front of you
  \item \emph{Cover:} write it again from memory, and count the pieces, not the strokes
\end{itemize}

\begin{tcolorbox}[breakable,colback=teal!4,colframe=teal!35!black,title={Before you move on}]
How many pieces in \zh{活}? (\textbf{Three}.) What is on the left? (\textbf{The water piece, as in \zh{洗}}.) What is under \zh{千}? (\textbf{\zh{口}}.)

Source: \href{https://raw.githubusercontent.com/chanind/hanzi-writer-data/68d10a4b21150cae5e1ebbd223eed289cf32d90c/data/\%E6\%B4\%BB.json}{Hanzi Writer Data \zh{活}, pinned stroke paths and medians}.
\end{tcolorbox}
\section[yīng'ér]{yīng'ér — newborn; infant}

\label{lesson:ZH-C302-yinger}



\begin{quote}
Before the new one: say the Mandarin for to be expecting a baby, then the Mandarin for to stay alive.
\end{quote}

\subsection*{You'll want to know: yīng'ér}
\emph{yīng'ér} — \textquotedblleft{}newborn; infant\textquotedblright{}.

\begin{quote}
\emph{zhège yīng'ér hěn kě'ài} — This baby is very sweet.
\end{quote}

Its characters come later; for now it is a word for the ear.

\begin{cousinweb}
\emph{yīng'ér} is built from \emph{yīng} and \zh{儿} 'a child', both inherited from Middle Chinese, \emph{'jieng} and \emph{nye} (Guangyun, in Baxter's transcription), which quotes an older word list for the pair: a baby girl is \emph{'jieng}, a baby boy \emph{nye}.
\end{cousinweb}

\begin{grammarlens}[title={What you've built}]
One more word for telling what happened.
\end{grammarlens}

\subsection*{Your turn}
\begin{itemize}
\raggedright
  \item \emph{Say it:} \emph{yīng'ér}
  \item \emph{Say it:} \emph{yīng'ér}, once more
  \item \emph{Say it:} \emph{zhège yīng'ér hěn kě'ài} — This baby is very sweet.
  \item \emph{Recall:} say the Mandarin for to be expecting a baby, then the Mandarin for to stay alive, then say \emph{yīng'ér} again
\end{itemize}

\begin{tcolorbox}[breakable,colback=teal!4,colframe=teal!35!black,title={Before you move on}]
What does yīng'ér mean? (\textquotedblleft{}Newborn; infant\textquotedblright{}.) Say it once more.
\end{tcolorbox}
\section[\texorpdfstring{\zh{救活} (jiùhuó)}{救活 (jiùhuó)}]{\zh{救活} (jiùhuó) — to bring back to life}

\label{lesson:ZH-C302-jiuhuo-revive}



\begin{quote}
Before the new one: say the Mandarin for to stay alive, then the Mandarin for a newborn.
\end{quote}

\subsection*{You'll want to know: \zh{救活}}
\textbf{\zh{救活}} — \emph{jiùhuó} — \textquotedblleft{}to bring back to life\textquotedblright{}.

\begin{quote}
\textbf{\zh{医生救活了他}} \emph{yīshēng jiùhuó le tā} — The doctor brought him back to life.
\end{quote}

\begin{cousinweb}
\textbf{\zh{救活}} (\emph{jiùhuó}) is built from \zh{救} 'to save' and \zh{活} 'alive', both inherited from Middle Chinese, \emph{kjuwH} and \emph{hwat} (Guangyun, in Baxter's transcription): to save and make alive.
\end{cousinweb}

\begin{grammarlens}[title={What you've built}]
One more word for telling what happened.
\end{grammarlens}

\subsection*{Your turn}
\begin{itemize}
\raggedright
  \item \emph{Say it:} \emph{jiùhuó}
  \item \emph{Say it:} \emph{jiùhuó}, once more
  \item \emph{Say it:} \emph{yīshēng jiùhuó le tā} — The doctor brought him back to life.
  \item \emph{Recall:} say the Mandarin for to stay alive, then the Mandarin for a newborn, then say \emph{jiùhuó} again
\end{itemize}

\begin{tcolorbox}[breakable,colback=teal!4,colframe=teal!35!black,title={Before you move on}]
What does \zh{救活} mean? (\textquotedblleft{}To bring back to life\textquotedblright{}.) Say it once more.
\end{tcolorbox}
\section[shuāngbāotāi]{shuāngbāotāi — twins}

\label{lesson:ZH-C302-shuangbaotai}



\begin{quote}
Before the new one: say the Mandarin for a newborn, then the Mandarin for to bring back to life.
\end{quote}

\subsection*{You'll want to know: shuāngbāotāi}
\emph{shuāngbāotāi} — \textquotedblleft{}twins\textquotedblright{}.

\begin{quote}
\emph{tāmen shì shuāngbāotāi} — They're twins.
\end{quote}

Its characters come later; for now it is a word for the ear.

\begin{cousinweb}
\emph{shuāngbāotāi} is built from \emph{shuāng} 'a pair', \emph{bāo} 'the womb' and \emph{tāi} 'an unborn child', all inherited from Middle Chinese, \emph{sraewng}, \emph{paew} and \emph{thoj} (Guangyun, in Baxter's transcription), which glosses the first 'two of a kind' and the third, quoting an older dictionary, 'three months with child', and gives the second only in the pair \emph{bāotāi}, so its sense 'the womb' is not spelled out there: two from one womb.
\end{cousinweb}

\begin{grammarlens}[title={What you've built}]
One more word for telling what happened.
\end{grammarlens}

\subsection*{Your turn}
\begin{itemize}
\raggedright
  \item \emph{Say it:} \emph{shuāngbāotāi}
  \item \emph{Say it:} \emph{shuāngbāotāi}, once more
  \item \emph{Say it:} \emph{tāmen shì shuāngbāotāi} — They're twins.
  \item \emph{Recall:} say the Mandarin for a newborn, then the Mandarin for to bring back to life, then say \emph{shuāngbāotāi} again
\end{itemize}

\begin{tcolorbox}[breakable,colback=teal!4,colframe=teal!35!black,title={Before you move on}]
What does shuāngbāotāi mean? (\textquotedblleft{}Twins\textquotedblright{}.) Say it once more.
\end{tcolorbox}
